% =============================================================================
%  progress.tex — รายงานความก้าวหน้าโครงการ ULMs
%
%  build:  cd ~/Projects/SoftwareEn/docs && latexmk progress.tex
%          ผลลัพธ์ออกที่ build/progress.pdf (ตาม $out_dir ใน .latexmkrc)
%
%  โครงเรื่องยึดตามสไลด์ Week 06 — PM: Project Progress Report ของอาจารย์
%  (resource/Week 06 - PM - Project Progress Report 2569-07-23 R04.pdf)
%  ตัวเลขความก้าวหน้าอ้างอิงหลักฐานจริงใน repo ดูที่มาในภาคผนวก ก
%
%  เขียนรายงานครั้งถัดไป: แก้ค่าใน "ตัวแปรของรอบรายงาน" ด้านล่าง แล้วไล่แก้
%  §3 (ตารางมิติงาน) · §4 (จุดบนกราฟ) · §5 (สถานะ Sprint) · §9–§11
% =============================================================================
% =============================================================================
%  ulms-preamble.tex — preamble ที่ใช้ร่วมกันทุกเอกสารของโครงการ ULMs
%
%  ใช้โดย:  main.tex (ข้อเสนอโครงการ) · progress.tex (รายงานความก้าวหน้า)
%
%  เอกสารที่เรียกใช้ต้องทำสองอย่างหลัง % =============================================================================
%  ulms-preamble.tex — preamble ที่ใช้ร่วมกันทุกเอกสารของโครงการ ULMs
%
%  ใช้โดย:  main.tex (ข้อเสนอโครงการ) · progress.tex (รายงานความก้าวหน้า)
%
%  เอกสารที่เรียกใช้ต้องทำสองอย่างหลัง % =============================================================================
%  ulms-preamble.tex — preamble ที่ใช้ร่วมกันทุกเอกสารของโครงการ ULMs
%
%  ใช้โดย:  main.tex (ข้อเสนอโครงการ) · progress.tex (รายงานความก้าวหน้า)
%
%  เอกสารที่เรียกใช้ต้องทำสองอย่างหลัง \input{ulms-preamble}:
%    1) \hypersetup{pdftitle={...}}  — ชื่อเอกสารใน metadata ของ PDF
%    2) \begin{document}
%  (pdftitle ถูกย้ายออกจากไฟล์นี้ เพราะเป็นค่าเฉพาะของแต่ละเอกสาร)
%
%  แก้ที่นี่ที่เดียว = มีผลกับทุกเอกสาร ห้ามก๊อป preamble ไปวางซ้ำ
% =============================================================================
% =============================================================================
%  main.tex  —  เอกสารข้อเสนอโครงการ (Academic Project Proposal)
%  ระบบบริหารจัดการการยืม–คืนอุปกรณ์มหาวิทยาลัย (ULMs)
%
%  IMPORTANT: This document contains Thai text and MUST be compiled with
%  XeLaTeX (not pdflatex). Thai needs a Unicode font + ICU line-breaking.
%      Build:  latexmk main.tex        (uses .latexmkrc -> xelatex)
%      or:     xelatex main.tex        (run twice for the table of contents)
% =============================================================================

\documentclass[11pt, a4paper]{article}

% ---- Fonts & Thai support ----------------------------------------------------
% fontspec lets XeLaTeX use any system font. Laksaman is a complete serif face
% that covers BOTH Thai and Latin glyphs, so one \setmainfont handles the whole
% (Thai + English) document without font switching.
\usepackage{fontspec}

% Thai has no spaces between words, so TeX cannot find line-break points on its
% own. These two XeTeX primitives switch on ICU's dictionary-based Thai word
% segmentation, which is what makes Thai paragraphs wrap correctly.
\XeTeXlinebreaklocale "th"
\XeTeXlinebreakskip = 0pt plus 0.1pt

\setmainfont{Laksaman}

% ---- Page geometry & paragraph style ----------------------------------------
\usepackage[margin=2.5cm]{geometry}
\usepackage{parskip}          % block paragraphs (no indent, gap between) — common
                              % and more readable for Thai documents
\linespread{1.15}             % a little extra leading; Thai has tall tone marks

% ---- Colours (monochrome / grayscale palette) ------------------------------
\usepackage[table]{xcolor}    % 'table' option enables \rowcolor for table headers
% Monochrome (grayscale) palette. Names kept as-is so every reference still
% works; only the tones changed from the original green/blue accents.
\definecolor{kugreen}{HTML}{1A1A1A}   % near-black — section headings & title
\definecolor{kudark}{HTML}{333333}    % dark gray — subsection/subsubsection titles
\definecolor{headbg}{HTML}{EAEAEA}    % light gray background for table header rows
\definecolor{linkblue}{HTML}{404040}  % medium-dark gray — links / citations

% ---- Section heading styling -------------------------------------------------
% 'nobottomtitles*' stops a heading from being printed near the bottom of a
% page: if less than \bottomtitlespace of room is left, the heading (and its
% following text) moves to the next page. This is titlesec's built-in, robust
% replacement for the earlier \needspace hack — the latter injected vertical
% glue that clashed with longtable page-splitting ("Infinite glue shrinkage").
\usepackage[nobottomtitles*]{titlesec}
\titleformat{\section}
  {\color{kugreen}\Large\bfseries}{\thesection}{0.6em}{}
\titleformat{\subsection}
  {\color{kudark}\large\bfseries}{\thesubsection}{0.5em}{}
\titleformat{\subsubsection}
  {\color{kudark}\normalsize\bfseries}{\thesubsubsection}{0.5em}{}
\titlespacing*{\section}{0pt}{1.4\baselineskip}{0.5\baselineskip}

% Minimum room a heading needs at the bottom of a page; if only ~1–3 lines are
% left (less than this), the heading is bumped to the next page instead of being
% stranded above its table/text.
\renewcommand{\bottomtitlespace}{4\baselineskip}

% ---- Hierarchical indentation by heading level ------------------------------
% Each heading level shifts BOTH the heading and the body text under it further
% to the right: section = 0, subsection (x.y) = 1 step, subsubsection (x.y.z) =
% 2 steps. We do this by having the sectioning commands set \leftskip (a global
% left margin that carries through every following line until the next heading
% resets it). \par first, so the *previous* paragraph is not re-indented.
% \leftskip indents plain paragraphs; \@totalleftmargin makes list environments
% (itemize/enumerate) start from the same indent instead of the page margin.
\makeatletter
\newlength{\lvlindent}
\setlength{\lvlindent}{1.8em}          % one indent step
\let\ulmsSecIndent\section
\renewcommand{\section}{\par\global\leftskip=0pt\global\@totalleftmargin=0pt\relax\ulmsSecIndent}
\let\ulmsSubsecIndent\subsection
\renewcommand{\subsection}{\par\global\leftskip=\lvlindent\global\@totalleftmargin=\lvlindent\relax\ulmsSubsecIndent}
\let\ulmsSubsubsecIndent\subsubsection
\renewcommand{\subsubsection}{\par\global\leftskip=2\lvlindent\global\@totalleftmargin=2\lvlindent\relax\ulmsSubsubsecIndent}
\makeatother

% Indent the headings themselves to match (titlesec ignores \leftskip for the
% heading line, so its left margin is set here via the first \titlespacing arg).
\titlespacing*{\subsection}   {\lvlindent} {1.0\baselineskip}{0.4\baselineskip}
\titlespacing*{\subsubsection}{2\lvlindent}{0.8\baselineskip}{0.3\baselineskip}

% ---- Lists (tighter, custom bullets) ----------------------------------------
\usepackage{enumitem}
\setlist[itemize]{leftmargin=1.6em, itemsep=2pt, topsep=3pt}
\setlist[enumerate]{leftmargin=1.9em, itemsep=2pt, topsep=3pt}

% ---- Tables ------------------------------------------------------------------
\usepackage{array}
\usepackage{booktabs}         % \toprule / \midrule / \bottomrule
\usepackage{tabularx}         % auto-width columns that wrap long Thai text
\usepackage{longtable}        % tables that can break across pages
\usepackage{multirow}         % merged cells for the team-structure table

% Left-aligned wrapping column types:
%   L{width} — fixed-width ragged-right paragraph column
%   Y        — flexible ragged-right X column (for tabularx)
\newcolumntype{L}[1]{>{\raggedright\arraybackslash}p{#1}}
\newcolumntype{Y}{>{\raggedright\arraybackslash}X}
\renewcommand{\arraystretch}{1.3}

% Helpers for a styled header row inside tables.
%   \hrow  — starts a header row and shades it (must be FIRST in the row;
%            colortbl allows only one \rowcolor per row).
%   \thead — bold header-cell text.
\newcommand{\hrow}{\rowcolor{headbg}}
\newcommand{\thead}[1]{\textbf{#1}}

% \begin{keeptable} … \end{keeptable} — keeps a table and the bold label above
% it on the SAME page (used by the data dictionary in §5). A tabularx cannot
% split across pages, so LaTeX is free to break between the label and its
% table, stranding the label alone at the bottom of a page; wrapping the pair
% in a minipage makes them one unbreakable block.
%   * \leftskip is zeroed inside: the minipage box is already placed at the
%     current heading indent, so leaving it set would indent the content twice
%     and push the table past the right margin.
%   * \parskip is restored because minipage's \@parboxrestore zeroes it, which
%     would glue the table straight onto the label line.
% Tables inside must use \linewidth (= the minipage width), not \textwidth.
\newenvironment{keeptable}
  {\par\medskip\noindent
   \begin{minipage}{\dimexpr\textwidth-\leftskip\relax}%
   \leftskip=0pt\relax
   \setlength{\parskip}{0.5\baselineskip plus 2pt}}
  {\end{minipage}\par\medskip}

% ---- Table of contents styling ----------------------------------------------
% By default LaTeX gives dotted leaders to subsections but NOT to sections
% (they print bold with just a right-aligned page number). tocloft lets us turn
% on dots for the section level too, so every TOC line runs "….." to its page.
\usepackage{tocloft}
\renewcommand{\cftsecdotsep}{\cftdotsep}   % dotted leaders for section entries
\renewcommand{\cftsecleader}{\cftdotfill{\cftsecdotsep}}
\renewcommand{\cftsecpagefont}{\normalfont}  % keep page numbers un-bold

% ---- Figures -----------------------------------------------------------------
\usepackage{graphicx}
\graphicspath{{figures/}}

% Landscape (rotated) full-page floats — needed by the ER diagram in §5, which
% is far wider than it is tall and would be unreadable scaled to \textwidth.
\usepackage{rotating}

% ---- Flowcharts / diagrams (TikZ) -------------------------------------------
% Used for the borrow–return workflow diagram (§5). Libraries: geometric shapes
% for decision diamonds, arrows.meta for arrowheads, positioning for relative
% node placement, calc/fit/backgrounds for routing loops and grouping.
\usepackage{tikz}
\usetikzlibrary{shapes.geometric, arrows.meta, positioning, calc, fit, backgrounds}

% ---- Links & bookmarks -------------------------------------------------------
\usepackage{hyperref}
\hypersetup{
  colorlinks=true, linkcolor=black, citecolor=linkblue, urlcolor=linkblue,
  pdflang={th},
}

% Thai labels for auto-generated headings.
\renewcommand{\contentsname}{สารบัญ (Table of Contents)}
\renewcommand{\tablename}{ตาราง}
\renewcommand{\figurename}{รูปที่}

:
%    1) \hypersetup{pdftitle={...}}  — ชื่อเอกสารใน metadata ของ PDF
%    2) \begin{document}
%  (pdftitle ถูกย้ายออกจากไฟล์นี้ เพราะเป็นค่าเฉพาะของแต่ละเอกสาร)
%
%  แก้ที่นี่ที่เดียว = มีผลกับทุกเอกสาร ห้ามก๊อป preamble ไปวางซ้ำ
% =============================================================================
% =============================================================================
%  main.tex  —  เอกสารข้อเสนอโครงการ (Academic Project Proposal)
%  ระบบบริหารจัดการการยืม–คืนอุปกรณ์มหาวิทยาลัย (ULMs)
%
%  IMPORTANT: This document contains Thai text and MUST be compiled with
%  XeLaTeX (not pdflatex). Thai needs a Unicode font + ICU line-breaking.
%      Build:  latexmk main.tex        (uses .latexmkrc -> xelatex)
%      or:     xelatex main.tex        (run twice for the table of contents)
% =============================================================================

\documentclass[11pt, a4paper]{article}

% ---- Fonts & Thai support ----------------------------------------------------
% fontspec lets XeLaTeX use any system font. Laksaman is a complete serif face
% that covers BOTH Thai and Latin glyphs, so one \setmainfont handles the whole
% (Thai + English) document without font switching.
\usepackage{fontspec}

% Thai has no spaces between words, so TeX cannot find line-break points on its
% own. These two XeTeX primitives switch on ICU's dictionary-based Thai word
% segmentation, which is what makes Thai paragraphs wrap correctly.
\XeTeXlinebreaklocale "th"
\XeTeXlinebreakskip = 0pt plus 0.1pt

\setmainfont{Laksaman}

% ---- Page geometry & paragraph style ----------------------------------------
\usepackage[margin=2.5cm]{geometry}
\usepackage{parskip}          % block paragraphs (no indent, gap between) — common
                              % and more readable for Thai documents
\linespread{1.15}             % a little extra leading; Thai has tall tone marks

% ---- Colours (monochrome / grayscale palette) ------------------------------
\usepackage[table]{xcolor}    % 'table' option enables \rowcolor for table headers
% Monochrome (grayscale) palette. Names kept as-is so every reference still
% works; only the tones changed from the original green/blue accents.
\definecolor{kugreen}{HTML}{1A1A1A}   % near-black — section headings & title
\definecolor{kudark}{HTML}{333333}    % dark gray — subsection/subsubsection titles
\definecolor{headbg}{HTML}{EAEAEA}    % light gray background for table header rows
\definecolor{linkblue}{HTML}{404040}  % medium-dark gray — links / citations

% ---- Section heading styling -------------------------------------------------
% 'nobottomtitles*' stops a heading from being printed near the bottom of a
% page: if less than \bottomtitlespace of room is left, the heading (and its
% following text) moves to the next page. This is titlesec's built-in, robust
% replacement for the earlier \needspace hack — the latter injected vertical
% glue that clashed with longtable page-splitting ("Infinite glue shrinkage").
\usepackage[nobottomtitles*]{titlesec}
\titleformat{\section}
  {\color{kugreen}\Large\bfseries}{\thesection}{0.6em}{}
\titleformat{\subsection}
  {\color{kudark}\large\bfseries}{\thesubsection}{0.5em}{}
\titleformat{\subsubsection}
  {\color{kudark}\normalsize\bfseries}{\thesubsubsection}{0.5em}{}
\titlespacing*{\section}{0pt}{1.4\baselineskip}{0.5\baselineskip}

% Minimum room a heading needs at the bottom of a page; if only ~1–3 lines are
% left (less than this), the heading is bumped to the next page instead of being
% stranded above its table/text.
\renewcommand{\bottomtitlespace}{4\baselineskip}

% ---- Hierarchical indentation by heading level ------------------------------
% Each heading level shifts BOTH the heading and the body text under it further
% to the right: section = 0, subsection (x.y) = 1 step, subsubsection (x.y.z) =
% 2 steps. We do this by having the sectioning commands set \leftskip (a global
% left margin that carries through every following line until the next heading
% resets it). \par first, so the *previous* paragraph is not re-indented.
% \leftskip indents plain paragraphs; \@totalleftmargin makes list environments
% (itemize/enumerate) start from the same indent instead of the page margin.
\makeatletter
\newlength{\lvlindent}
\setlength{\lvlindent}{1.8em}          % one indent step
\let\ulmsSecIndent\section
\renewcommand{\section}{\par\global\leftskip=0pt\global\@totalleftmargin=0pt\relax\ulmsSecIndent}
\let\ulmsSubsecIndent\subsection
\renewcommand{\subsection}{\par\global\leftskip=\lvlindent\global\@totalleftmargin=\lvlindent\relax\ulmsSubsecIndent}
\let\ulmsSubsubsecIndent\subsubsection
\renewcommand{\subsubsection}{\par\global\leftskip=2\lvlindent\global\@totalleftmargin=2\lvlindent\relax\ulmsSubsubsecIndent}
\makeatother

% Indent the headings themselves to match (titlesec ignores \leftskip for the
% heading line, so its left margin is set here via the first \titlespacing arg).
\titlespacing*{\subsection}   {\lvlindent} {1.0\baselineskip}{0.4\baselineskip}
\titlespacing*{\subsubsection}{2\lvlindent}{0.8\baselineskip}{0.3\baselineskip}

% ---- Lists (tighter, custom bullets) ----------------------------------------
\usepackage{enumitem}
\setlist[itemize]{leftmargin=1.6em, itemsep=2pt, topsep=3pt}
\setlist[enumerate]{leftmargin=1.9em, itemsep=2pt, topsep=3pt}

% ---- Tables ------------------------------------------------------------------
\usepackage{array}
\usepackage{booktabs}         % \toprule / \midrule / \bottomrule
\usepackage{tabularx}         % auto-width columns that wrap long Thai text
\usepackage{longtable}        % tables that can break across pages
\usepackage{multirow}         % merged cells for the team-structure table

% Left-aligned wrapping column types:
%   L{width} — fixed-width ragged-right paragraph column
%   Y        — flexible ragged-right X column (for tabularx)
\newcolumntype{L}[1]{>{\raggedright\arraybackslash}p{#1}}
\newcolumntype{Y}{>{\raggedright\arraybackslash}X}
\renewcommand{\arraystretch}{1.3}

% Helpers for a styled header row inside tables.
%   \hrow  — starts a header row and shades it (must be FIRST in the row;
%            colortbl allows only one \rowcolor per row).
%   \thead — bold header-cell text.
\newcommand{\hrow}{\rowcolor{headbg}}
\newcommand{\thead}[1]{\textbf{#1}}

% \begin{keeptable} … \end{keeptable} — keeps a table and the bold label above
% it on the SAME page (used by the data dictionary in §5). A tabularx cannot
% split across pages, so LaTeX is free to break between the label and its
% table, stranding the label alone at the bottom of a page; wrapping the pair
% in a minipage makes them one unbreakable block.
%   * \leftskip is zeroed inside: the minipage box is already placed at the
%     current heading indent, so leaving it set would indent the content twice
%     and push the table past the right margin.
%   * \parskip is restored because minipage's \@parboxrestore zeroes it, which
%     would glue the table straight onto the label line.
% Tables inside must use \linewidth (= the minipage width), not \textwidth.
\newenvironment{keeptable}
  {\par\medskip\noindent
   \begin{minipage}{\dimexpr\textwidth-\leftskip\relax}%
   \leftskip=0pt\relax
   \setlength{\parskip}{0.5\baselineskip plus 2pt}}
  {\end{minipage}\par\medskip}

% ---- Table of contents styling ----------------------------------------------
% By default LaTeX gives dotted leaders to subsections but NOT to sections
% (they print bold with just a right-aligned page number). tocloft lets us turn
% on dots for the section level too, so every TOC line runs "….." to its page.
\usepackage{tocloft}
\renewcommand{\cftsecdotsep}{\cftdotsep}   % dotted leaders for section entries
\renewcommand{\cftsecleader}{\cftdotfill{\cftsecdotsep}}
\renewcommand{\cftsecpagefont}{\normalfont}  % keep page numbers un-bold

% ---- Figures -----------------------------------------------------------------
\usepackage{graphicx}
\graphicspath{{figures/}}

% Landscape (rotated) full-page floats — needed by the ER diagram in §5, which
% is far wider than it is tall and would be unreadable scaled to \textwidth.
\usepackage{rotating}

% ---- Flowcharts / diagrams (TikZ) -------------------------------------------
% Used for the borrow–return workflow diagram (§5). Libraries: geometric shapes
% for decision diamonds, arrows.meta for arrowheads, positioning for relative
% node placement, calc/fit/backgrounds for routing loops and grouping.
\usepackage{tikz}
\usetikzlibrary{shapes.geometric, arrows.meta, positioning, calc, fit, backgrounds}

% ---- Links & bookmarks -------------------------------------------------------
\usepackage{hyperref}
\hypersetup{
  colorlinks=true, linkcolor=black, citecolor=linkblue, urlcolor=linkblue,
  pdflang={th},
}

% Thai labels for auto-generated headings.
\renewcommand{\contentsname}{สารบัญ (Table of Contents)}
\renewcommand{\tablename}{ตาราง}
\renewcommand{\figurename}{รูปที่}

:
%    1) \hypersetup{pdftitle={...}}  — ชื่อเอกสารใน metadata ของ PDF
%    2) \begin{document}
%  (pdftitle ถูกย้ายออกจากไฟล์นี้ เพราะเป็นค่าเฉพาะของแต่ละเอกสาร)
%
%  แก้ที่นี่ที่เดียว = มีผลกับทุกเอกสาร ห้ามก๊อป preamble ไปวางซ้ำ
% =============================================================================
% =============================================================================
%  main.tex  —  เอกสารข้อเสนอโครงการ (Academic Project Proposal)
%  ระบบบริหารจัดการการยืม–คืนอุปกรณ์มหาวิทยาลัย (ULMs)
%
%  IMPORTANT: This document contains Thai text and MUST be compiled with
%  XeLaTeX (not pdflatex). Thai needs a Unicode font + ICU line-breaking.
%      Build:  latexmk main.tex        (uses .latexmkrc -> xelatex)
%      or:     xelatex main.tex        (run twice for the table of contents)
% =============================================================================

\documentclass[11pt, a4paper]{article}

% ---- Fonts & Thai support ----------------------------------------------------
% fontspec lets XeLaTeX use any system font. Laksaman is a complete serif face
% that covers BOTH Thai and Latin glyphs, so one \setmainfont handles the whole
% (Thai + English) document without font switching.
\usepackage{fontspec}

% Thai has no spaces between words, so TeX cannot find line-break points on its
% own. These two XeTeX primitives switch on ICU's dictionary-based Thai word
% segmentation, which is what makes Thai paragraphs wrap correctly.
\XeTeXlinebreaklocale "th"
\XeTeXlinebreakskip = 0pt plus 0.1pt

\setmainfont{Laksaman}

% ---- Page geometry & paragraph style ----------------------------------------
\usepackage[margin=2.5cm]{geometry}
\usepackage{parskip}          % block paragraphs (no indent, gap between) — common
                              % and more readable for Thai documents
\linespread{1.15}             % a little extra leading; Thai has tall tone marks

% ---- Colours (monochrome / grayscale palette) ------------------------------
\usepackage[table]{xcolor}    % 'table' option enables \rowcolor for table headers
% Monochrome (grayscale) palette. Names kept as-is so every reference still
% works; only the tones changed from the original green/blue accents.
\definecolor{kugreen}{HTML}{1A1A1A}   % near-black — section headings & title
\definecolor{kudark}{HTML}{333333}    % dark gray — subsection/subsubsection titles
\definecolor{headbg}{HTML}{EAEAEA}    % light gray background for table header rows
\definecolor{linkblue}{HTML}{404040}  % medium-dark gray — links / citations

% ---- Section heading styling -------------------------------------------------
% 'nobottomtitles*' stops a heading from being printed near the bottom of a
% page: if less than \bottomtitlespace of room is left, the heading (and its
% following text) moves to the next page. This is titlesec's built-in, robust
% replacement for the earlier \needspace hack — the latter injected vertical
% glue that clashed with longtable page-splitting ("Infinite glue shrinkage").
\usepackage[nobottomtitles*]{titlesec}
\titleformat{\section}
  {\color{kugreen}\Large\bfseries}{\thesection}{0.6em}{}
\titleformat{\subsection}
  {\color{kudark}\large\bfseries}{\thesubsection}{0.5em}{}
\titleformat{\subsubsection}
  {\color{kudark}\normalsize\bfseries}{\thesubsubsection}{0.5em}{}
\titlespacing*{\section}{0pt}{1.4\baselineskip}{0.5\baselineskip}

% Minimum room a heading needs at the bottom of a page; if only ~1–3 lines are
% left (less than this), the heading is bumped to the next page instead of being
% stranded above its table/text.
\renewcommand{\bottomtitlespace}{4\baselineskip}

% ---- Hierarchical indentation by heading level ------------------------------
% Each heading level shifts BOTH the heading and the body text under it further
% to the right: section = 0, subsection (x.y) = 1 step, subsubsection (x.y.z) =
% 2 steps. We do this by having the sectioning commands set \leftskip (a global
% left margin that carries through every following line until the next heading
% resets it). \par first, so the *previous* paragraph is not re-indented.
% \leftskip indents plain paragraphs; \@totalleftmargin makes list environments
% (itemize/enumerate) start from the same indent instead of the page margin.
\makeatletter
\newlength{\lvlindent}
\setlength{\lvlindent}{1.8em}          % one indent step
\let\ulmsSecIndent\section
\renewcommand{\section}{\par\global\leftskip=0pt\global\@totalleftmargin=0pt\relax\ulmsSecIndent}
\let\ulmsSubsecIndent\subsection
\renewcommand{\subsection}{\par\global\leftskip=\lvlindent\global\@totalleftmargin=\lvlindent\relax\ulmsSubsecIndent}
\let\ulmsSubsubsecIndent\subsubsection
\renewcommand{\subsubsection}{\par\global\leftskip=2\lvlindent\global\@totalleftmargin=2\lvlindent\relax\ulmsSubsubsecIndent}
\makeatother

% Indent the headings themselves to match (titlesec ignores \leftskip for the
% heading line, so its left margin is set here via the first \titlespacing arg).
\titlespacing*{\subsection}   {\lvlindent} {1.0\baselineskip}{0.4\baselineskip}
\titlespacing*{\subsubsection}{2\lvlindent}{0.8\baselineskip}{0.3\baselineskip}

% ---- Lists (tighter, custom bullets) ----------------------------------------
\usepackage{enumitem}
\setlist[itemize]{leftmargin=1.6em, itemsep=2pt, topsep=3pt}
\setlist[enumerate]{leftmargin=1.9em, itemsep=2pt, topsep=3pt}

% ---- Tables ------------------------------------------------------------------
\usepackage{array}
\usepackage{booktabs}         % \toprule / \midrule / \bottomrule
\usepackage{tabularx}         % auto-width columns that wrap long Thai text
\usepackage{longtable}        % tables that can break across pages
\usepackage{multirow}         % merged cells for the team-structure table

% Left-aligned wrapping column types:
%   L{width} — fixed-width ragged-right paragraph column
%   Y        — flexible ragged-right X column (for tabularx)
\newcolumntype{L}[1]{>{\raggedright\arraybackslash}p{#1}}
\newcolumntype{Y}{>{\raggedright\arraybackslash}X}
\renewcommand{\arraystretch}{1.3}

% Helpers for a styled header row inside tables.
%   \hrow  — starts a header row and shades it (must be FIRST in the row;
%            colortbl allows only one \rowcolor per row).
%   \thead — bold header-cell text.
\newcommand{\hrow}{\rowcolor{headbg}}
\newcommand{\thead}[1]{\textbf{#1}}

% \begin{keeptable} … \end{keeptable} — keeps a table and the bold label above
% it on the SAME page (used by the data dictionary in §5). A tabularx cannot
% split across pages, so LaTeX is free to break between the label and its
% table, stranding the label alone at the bottom of a page; wrapping the pair
% in a minipage makes them one unbreakable block.
%   * \leftskip is zeroed inside: the minipage box is already placed at the
%     current heading indent, so leaving it set would indent the content twice
%     and push the table past the right margin.
%   * \parskip is restored because minipage's \@parboxrestore zeroes it, which
%     would glue the table straight onto the label line.
% Tables inside must use \linewidth (= the minipage width), not \textwidth.
\newenvironment{keeptable}
  {\par\medskip\noindent
   \begin{minipage}{\dimexpr\textwidth-\leftskip\relax}%
   \leftskip=0pt\relax
   \setlength{\parskip}{0.5\baselineskip plus 2pt}}
  {\end{minipage}\par\medskip}

% ---- Table of contents styling ----------------------------------------------
% By default LaTeX gives dotted leaders to subsections but NOT to sections
% (they print bold with just a right-aligned page number). tocloft lets us turn
% on dots for the section level too, so every TOC line runs "….." to its page.
\usepackage{tocloft}
\renewcommand{\cftsecdotsep}{\cftdotsep}   % dotted leaders for section entries
\renewcommand{\cftsecleader}{\cftdotfill{\cftsecdotsep}}
\renewcommand{\cftsecpagefont}{\normalfont}  % keep page numbers un-bold

% ---- Figures -----------------------------------------------------------------
\usepackage{graphicx}
\graphicspath{{figures/}}

% Landscape (rotated) full-page floats — needed by the ER diagram in §5, which
% is far wider than it is tall and would be unreadable scaled to \textwidth.
\usepackage{rotating}

% ---- Flowcharts / diagrams (TikZ) -------------------------------------------
% Used for the borrow–return workflow diagram (§5). Libraries: geometric shapes
% for decision diamonds, arrows.meta for arrowheads, positioning for relative
% node placement, calc/fit/backgrounds for routing loops and grouping.
\usepackage{tikz}
\usetikzlibrary{shapes.geometric, arrows.meta, positioning, calc, fit, backgrounds}

% ---- Links & bookmarks -------------------------------------------------------
\usepackage{hyperref}
\hypersetup{
  colorlinks=true, linkcolor=black, citecolor=linkblue, urlcolor=linkblue,
  pdflang={th},
}

% Thai labels for auto-generated headings.
\renewcommand{\contentsname}{สารบัญ (Table of Contents)}
\renewcommand{\tablename}{ตาราง}
\renewcommand{\figurename}{รูปที่}



\hypersetup{pdftitle={รายงานความก้าวหน้าโครงการ ULMs ครั้งที่ 1}}

% ---- ตัวแปรของรอบรายงาน — แก้ที่นี่ที่เดียวเวลาออกรายงานรอบใหม่ ----------------
\newcommand{\reportno}{1}
\newcommand{\reportdate}{7 สิงหาคม 2569}
\newcommand{\reportperiod}{24 กรกฎาคม – 7 สิงหาคม 2569}
\newcommand{\plannedpct}{43\%}
\newcommand{\actualpct}{24\%}

\begin{document}

% =============================================================================
%  TITLE PAGE
% =============================================================================
\begin{titlepage}
\centering
\vspace*{2.5cm}

{\Large\bfseries\color{kugreen} รายงานความก้าวหน้าโครงการ ครั้งที่ \reportno\par}
\vspace{0.6cm}
{\LARGE\bfseries\color{kugreen} ระบบบริหารจัดการการยืม–คืนอุปกรณ์มหาวิทยาลัย\par}
\vspace{0.35cm}
{\Large\color{kudark} University Lending Management System (ULMs)\par}

\vspace{1.6cm}
\rule{0.72\textwidth}{0.8pt}
\vspace{0.9cm}

\begin{tabular}{@{}r@{\hspace{1.2em}}l@{}}
\textbf{รอบรายงาน}    & \reportperiod \\[0.35em]
\textbf{วันที่รายงาน}  & \reportdate \\[0.35em]
\textbf{สถานะ Sprint} & Sprint 1 ``โครงเดินได้'' (วันที่ 8 จาก 14) \\[0.35em]
\textbf{ความก้าวหน้า} & \actualpct\ เทียบแผน \plannedpct \\[0.35em]
\textbf{สรุปสถานะ}    & \textbf{ช้ากว่าแผน (Behind Schedule)} \\
\end{tabular}

\vspace{0.9cm}
\rule{0.72\textwidth}{0.8pt}

\vfill
{\large\color{kudark} คณะทำงาน ``miro ปั่น 14 บาท''\par}
\vspace{0.3cm}
{\color{kudark} ภาควิชาวิศวกรรมคอมพิวเตอร์ คณะวิศวกรรมศาสตร์\par}
{\color{kudark} มหาวิทยาลัยเกษตรศาสตร์\par}
\vspace{1.2cm}
\end{titlepage}

\tableofcontents
\clearpage

% =============================================================================
\section{บทสรุปผู้บริหาร (Executive Summary)}

ณ วันที่ \reportdate\ โครงการอยู่ระหว่าง Sprint 1 ``โครงเดินได้'' เป็นวันที่ 8 จากทั้งหมด 14 วัน คิดเป็นระยะเวลาที่ใช้ไปแล้ว 2 สัปดาห์จากแผนทั้งหมด 10 สัปดาห์

\textbf{ความก้าวหน้าโดยรวมอยู่ที่ \actualpct\ เทียบกับแผนที่ควรอยู่ที่ \plannedpct\ นั่นคือช้ากว่าแผนประมาณ 19 จุด} สาเหตุหลักไม่ได้มาจากปริมาณงานที่ทำได้น้อย แต่มาจาก\textbf{องค์ประกอบของงานที่ทำไม่ตรงกับเส้นทางวิกฤต (Critical Path)} กล่าวคือ งานฝั่งออกแบบฐานข้อมูลและงานเอกสารเดินหน้าไปได้ดีถึงเกินแผนเล็กน้อย ขณะที่งานที่เป็นเงื่อนไขให้ Sprint ถัดไปเดินต่อได้ยังไม่เสร็จ

\textbf{ประเด็นวิกฤตที่สุดคือสัญญา tRPC (tRPC contract) ซึ่งเป็นงานของ Sprint 0 ยังไม่ถูกจัดทำ} ส่งผลให้ทีม Front-End และ Back-End ต่างพัฒนาแยกกันบนข้อมูลจำลอง (Mock) คนละชุด และจนถึงวันที่รายงานนี้ \textbf{ระบบทั้งสองฝั่งยังไม่เคยเชื่อมต่อกันจริงแม้แต่ครั้งเดียว} เป้าหมายส่งมอบของ Sprint 1 ที่กำหนดว่าต้องยืมอุปกรณ์ข้ามภาควิชาได้ครบวงจร จึงมีความเสี่ยงสูงมากที่จะไม่สำเร็จภายในวันที่ 13 สิงหาคม 2569

ในด้านงบประมาณ โครงการยังไม่มีค่าใช้จ่ายจริงเกิดขึ้น เนื่องจากยังไม่ถึงขั้นนำระบบขึ้นให้บริการ คิดเป็นการใช้งบประมาณ 0\% จากวงเงินประเมิน 5{,}000 บาท

\subsection*{ข้อเสนอเร่งด่วนต่อคณะกรรมการ}

\begin{enumerate}
  \item ให้ถือว่า \textbf{การจัดทำสัญญา tRPC เป็นงานเร่งด่วนสูงสุด} และต้องเสร็จภายในวันที่ 9 สิงหาคม 2569 มิฉะนั้นให้ปรับเป้าหมายของ Sprint 1 ลง
  \item \textbf{ปรับเป้าหมายส่งมอบของ Sprint 1} จาก ``ยืมข้ามภาคครบวงจร'' เหลือ ``เชื่อม Front-End กับ Back-End ได้จริงหนึ่งเส้นทาง'' เพื่อพิสูจน์ว่าสถาปัตยกรรมต่อกันติดก่อน แล้วเลื่อนงานที่เหลือไป Sprint 2
  \item \textbf{เพิ่มกำลังคนฝั่ง Back-End} ซึ่งขณะนี้มีผู้ลงมือเขียนโค้ดจริงเพียงคนเดียว สวนทางกับสัดส่วนงานที่ฝั่ง Back-End รับผิดชอบสูงที่สุดในโครงการ
\end{enumerate}

% =============================================================================
\section{สถานะโครงการโดยรวม (เงิน · งาน · เวลา)}

\begin{tabularx}{\dimexpr\textwidth-\leftskip\relax}{@{} L{2.6cm} L{3.2cm} L{3.2cm} Y @{}}
\toprule
\hrow \thead{ด้าน} & \thead{แผน} & \thead{ผลจริง} & \thead{การประเมิน} \\
\midrule
เวลา & ผ่านไป 2 จาก 10 สัปดาห์ (20\%) & ผ่านไป 2 จาก 10 สัปดาห์ (20\%) & เป็นไปตามปฏิทิน ยังไม่มีการเลื่อนวันสิ้นสุดโครงการ \\
งาน & \plannedpct & \actualpct & \textbf{ช้ากว่าแผน 19 จุด} กระจุกอยู่ที่งานพัฒนา Back-End และงานทดสอบ \\
เงิน & 5{,}000 บาท (วงเงินประเมินทั้งโครงการ) & 0 บาท (0\%) & ยังไม่มีรายจ่าย เพราะยังไม่ถึงขั้นเช่า Cloud/Hosting ตามที่ประเมินไว้ในข้อเสนอโครงการ \\
\bottomrule
\end{tabularx}

\vspace{0.8em}
\noindent ข้อสังเกตสำคัญคือ \textbf{ค่าใช้จ่าย 0\% ไม่ใช่สัญญาณที่ดีในกรณีนี้} เพราะงบประมาณที่ตั้งไว้ผูกกับการนำระบบขึ้นใช้งานจริง (Cloud + Database 1{,}500 บาท · Hosting 1{,}200 บาท · Domain 500 บาท) การที่ยังไม่มีรายจ่ายเลยสะท้อนว่างานยังไม่เดินไปถึงขั้นนั้น ไม่ได้แปลว่าประหยัดงบได้

% =============================================================================
\section{ความก้าวหน้าแยกตามมิติของงาน}
\label{sec:dimension}

การแบ่งมิติของงานในตารางนี้อ้างอิงแนวทางที่กำหนดไว้ในหัวข้อ Weekly Progress Report โดยแยกเป็นส่วนออกแบบ ส่วนพัฒนา ส่วนทดสอบ และส่วนเอกสาร ค่าน้ำหนักสะท้อนสัดส่วนปริมาณงานของแต่ละมิติในโครงการทั้งหมด และตัวเลขผลจริงประเมินจากหลักฐานที่ตรวจสอบได้ใน Repository ซึ่งแจกแจงไว้ในภาคผนวก ก

\begin{tabularx}{\dimexpr\textwidth-\leftskip\relax}{@{} L{4.4cm} L{1.3cm} L{1.2cm} L{1.2cm} L{1.5cm} Y @{}}
\toprule
\hrow \thead{มิติงาน} & \thead{น้ำหนัก} & \thead{แผน} & \thead{จริง} & \thead{ผลต่าง} & \thead{หมายเหตุ} \\
\midrule
Back-End Design \newline (API + Database Model) & 10\% & 90\% & 45\% & $-45$ & แบบจำลองข้อมูลเสร็จ แต่สัญญา API ยังไม่เริ่ม \\
Front-End Design \newline (UI/UX + Business Process) & 10\% & 65\% & 60\% & $-5$ & กระบวนการทางธุรกิจครบ แต่ยังไม่มี Wireframe ของหน้าจอส่วนใหญ่ \\
Back-End Dev & 30\% & 35\% & 12\% & $\mathbf{-23}$ & มีเฉพาะโครงสร้างพื้นฐาน ยังไม่มีฟังก์ชันงานจริง \\
Front-End Dev & 25\% & 40\% & 20\% & $\mathbf{-20}$ & มีเปลือกระบบและหน้าเข้าสู่ระบบ ที่เหลือเป็นหน้าเปล่า \\
Test (Front-End + Back-End) & 15\% & 20\% & 0\% & $-20$ & ยังไม่มีแผนทดสอบและยังไม่มีชุดทดสอบใด \\
Documentation & 10\% & 40\% & 45\% & $+5$ & \textbf{เร็วกว่าแผน} ข้อเสนอโครงการ 42 หน้าเสร็จสมบูรณ์ \\
\midrule
\hrow \thead{รวมถ่วงน้ำหนัก} & \thead{100\%} & \thead{\plannedpct} & \thead{\actualpct} & \thead{$\mathbf{-19}$} & \\
\bottomrule
\end{tabularx}

\vspace{0.8em}
\noindent อ่านตารางนี้แล้วจะเห็นรูปแบบที่ชัดเจน คือ \textbf{มิติที่ทำได้ตามแผนหรือดีกว่าแผนล้วนเป็นงานที่ทำคนเดียวได้และไม่ต้องรอใคร} ได้แก่ เอกสารและการออกแบบฐานข้อมูล ส่วนมิติที่ตกแผนหนักล้วนเป็นงานที่ต้องอาศัยข้อตกลงร่วมกันระหว่างสองฝั่งก่อนจึงจะเริ่มได้ ซึ่งชี้กลับไปที่ต้นเหตุเดียวกันคือสัญญา API ที่ยังไม่ถูกจัดทำ

% =============================================================================
\section{กราฟความก้าวหน้าเทียบแผน (S-Curve)}

\begin{figure}[h]
\centering
\begin{tikzpicture}[x=1.05cm, y=0.055cm]
  % ---- กริดและแกน ----
  \foreach \y in {0,20,40,60,80,100}
    \draw[gray!25] (0,\y) -- (10.4,\y);
  \draw[->, thick] (0,0) -- (10.8,0) node[right, font=\small] {สัปดาห์};
  \draw[->, thick] (0,0) -- (0,112) node[above, font=\small] {\% ความก้าวหน้า};
  \foreach \y in {0,20,40,60,80,100}
    \draw (0,\y) -- (-0.12,\y) node[left, font=\scriptsize] {\y};
  \foreach \x/\lbl in {0/{24 ก.ค.}, 2/{7 ส.ค.}, 4/{21 ส.ค.}, 6/{4 ก.ย.}, 8/{18 ก.ย.}, 10/{1 ต.ค.}}
    \draw (\x,0) -- (\x,-2) node[below, font=\scriptsize] {\lbl};

  % ---- เส้นแผน (Planned) ----
  \draw[thick, kudark, dashed]
    plot coordinates {(0,0) (1,25) (3,55) (5,75) (7,90) (9,97) (10,100)};
  \foreach \p in {(0,0), (1,25), (3,55), (5,75), (7,90), (9,97), (10,100)}
    \fill[kudark] \p circle (1.6pt);

  % ---- เส้นผลจริง (Actual) ----
  \draw[very thick, kugreen]
    plot coordinates {(0,0) (1,14) (2,24)};
  \foreach \p in {(0,0), (1,14), (2,24)}
    \fill[kugreen] \p circle (2.2pt);

  % ---- ช่องว่างระหว่างแผนกับผลจริง ณ วันที่รายงาน ----
  \draw[<->, red!70!black, thick] (2,24) -- (2,43);
  \node[right, font=\scriptsize, red!70!black] at (2.1,33.5) {ช้ากว่าแผน 19 จุด};
  \fill[red!70!black] (2,43) circle (1.6pt);

  % ---- คำอธิบายสัญลักษณ์ ----
  \draw[thick, kudark, dashed] (5.6,18) -- (6.4,18);
  \node[right, font=\scriptsize] at (6.5,18) {แผน (Planned)};
  \draw[very thick, kugreen] (5.6,8) -- (6.4,8);
  \node[right, font=\scriptsize] at (6.5,8) {ผลจริง (Actual)};
\end{tikzpicture}
\caption{กราฟ S-Curve เปรียบเทียบความก้าวหน้าตามแผนกับผลการดำเนินงานจริง
  เส้นผลจริงสิ้นสุดที่สัปดาห์ที่ 2 ซึ่งเป็นวันที่ออกรายงานฉบับนี้
  ค่าความก้าวหน้าคำนวณแบบถ่วงน้ำหนักตามตารางในหัวข้อ \ref{sec:dimension}}
\label{fig:scurve}
\end{figure}

\noindent รูปที่ \ref{fig:scurve} แสดงให้เห็นว่าช่องว่างระหว่างแผนกับผลจริงเริ่มถ่างออกตั้งแต่ปลาย Sprint 0 และกว้างขึ้นในสัปดาห์ที่ 2 หากแนวโน้มนี้ดำเนินต่อไปโดยไม่มีการแก้ไข ช่องว่างจะยิ่งขยายในช่วง Sprint 2 ถึง Sprint 3 ซึ่งเป็นช่วงที่แผนกำหนดปริมาณงานไว้หนาแน่นที่สุด

% =============================================================================
\clearpage
\section{แผนและผลการดำเนินงานที่ผ่านมา}

\subsection{Sprint 0 ``วางฐาน'' (24 – 30 กรกฎาคม 2569)}

\noindent\textbf{เป้าหมายตามแผน} ล็อกสัญญา API และโครงสร้างฐานข้อมูล โดยรวมความสามารถด้านหลายหน่วยงาน (Multi-unit) และการจองล่วงหน้าไว้ตั้งแต่ต้น เนื่องจากการเพิ่มความสามารถทั้งสองอย่างในภายหลังจะต้องรื้อระบบใหม่ทั้งหมด

\begin{tabularx}{\dimexpr\textwidth-\leftskip\relax}{@{} L{1.2cm} Y L{2.3cm} @{}}
\toprule
\hrow \thead{ส่วนงาน} & \thead{งานตามแผน} & \thead{สถานะจริง} \\
\midrule
SA & ปิดขอบเขต: multi-unit + facility เป็นงานหลัก · การจองตัดได้ · consumable/อุทธรณ์ เลื่อนออก & \textbf{เสร็จ} \\
BE & Prisma schema แกนกลาง รวมหน่วยงานผู้ให้ยืม สมาชิกตามหน่วยงาน เจ้าของทรัพยากร การจอง และการผูกชิ้น & \textbf{เสร็จบางส่วน} \\
BE & สัญญา tRPC ของวงจรหลัก แบบอิงสมาชิกตามหน่วยงาน & \textbf{ยังไม่เริ่ม} \\
BE & ติดตั้ง NestJS + tRPC + Prisma + PostgreSQL พร้อม migration และข้อมูลตั้งต้น 2 หน่วยงาน & \textbf{เสร็จบางส่วน} \\
FE & ติดตั้ง Vite + React + TypeScript + Router พร้อม CI และโครงการยืนยันตัวตน & \textbf{เสร็จบางส่วน} \\
FE & ระบบออกแบบ (Tailwind + shadcn) และ Wireframe หน้าเข้าสู่ระบบกับเปลือกระบบ & \textbf{เสร็จ} \\
FE & ชั้นข้อมูลจำลองที่สร้างจากสัญญา API & \textbf{ยังไม่เริ่ม} \\
QA & แผนการทดสอบและเกณฑ์การยอมรับ & \textbf{ยังไม่เริ่ม} \\
PM & ตั้งกระดานงาน พิธีการประจำ Sprint และเกณฑ์ ``เสร็จ'' & \textbf{ยังไม่เริ่ม} \\
\bottomrule
\end{tabularx}

\vspace{0.8em}
\noindent\textbf{เกณฑ์ ``เสร็จ'' ของ Sprint 0} กำหนดไว้ว่า สัญญาและโครงสร้างฐานข้อมูลต้องถูกล็อก มีข้อมูลตั้งต้น 2 หน่วยงาน และ Repository ต้องรันได้ทั้งฝั่ง Front-End และ Back-End \textbf{ผลคือยังไม่ผ่านเกณฑ์} เพราะขาดสัญญา API และขาดข้อมูลตั้งต้น ฝั่ง Back-End จึงยังรันเพื่อทดสอบวงจรงานจริงไม่ได้

\subsection{Sprint 1 ``โครงเดินได้'' (31 กรกฎาคม – 13 สิงหาคม 2569)}

\noindent\textbf{เป้าหมายตามแผน} ทำให้เกิดเส้นทางการทำงานที่ทะลุทุกชั้นของระบบ (Vertical Slice) โดยยืมอุปกรณ์ข้ามหน่วยงานได้ในเส้นทางปกติ เพื่อพิสูจน์ว่าสถาปัตยกรรมและกลไกสมาชิกตามหน่วยงาน ต่อกันได้จริง

\begin{tabularx}{\dimexpr\textwidth-\leftskip\relax}{@{} L{1.2cm} Y L{2.3cm} @{}}
\toprule
\hrow \thead{ส่วนงาน} & \thead{งานตามแผน} & \thead{สถานะจริง} \\
\midrule
BE & ยืนยันตัวตนและควบคุมสิทธิ์แบบอิงสมาชิกต่อหน่วยงาน & \textbf{ยังไม่เริ่ม} \\
BE & รายการอุปกรณ์ ค้นหา จำนวนคงเหลือ และการแสดงหน่วยงานเจ้าของ & \textbf{ยังไม่เริ่ม} \\
BE & สร้างคำขอ อนุมัติอัตโนมัติ ผูกชิ้น ส่งมอบ และรับคืน & \textbf{ยังไม่เริ่ม} \\
FE & หน้าเข้าสู่ระบบและการคุ้มกันเส้นทางตามสมาชิกตามหน่วยงาน & \textbf{เสร็จบางส่วน} \\
FE & หน้าเรียกดูอุปกรณ์ รายละเอียด และจำนวนคงเหลือ & \textbf{ยังไม่เริ่ม} \\
FE & หน้าส่งคำขอและติดตามสถานะของผู้ยืม & \textbf{ยังไม่เริ่ม} \\
FE & หน้าคิวงานของเจ้าหน้าที่ & \textbf{ยังไม่เริ่ม} \\
QA & ทดสอบเส้นทางปกติและกรณียืมข้ามภาควิชา & \textbf{ยังไม่เริ่ม} \\
\bottomrule
\end{tabularx}

\vspace{0.8em}
\noindent ผ่านไปแล้ว 8 วันจาก 14 วันของ Sprint 1 แต่ยังไม่มีงานใดในรายการข้างต้นที่เสร็จสมบูรณ์ งานที่คืบหน้าเพียงรายการเดียวคือหน้าเข้าสู่ระบบ ซึ่งทำงานอยู่บนการยืนยันตัวตนจำลอง ยังไม่ได้เชื่อมกับระบบยืนยันตัวตนจริง

% =============================================================================
\clearpage
\section{ผลงานที่ส่งมอบแล้ว}

หัวข้อนี้รายงานเฉพาะผลงานที่ตรวจสอบได้จริงใน Repository ไม่รวมกิจกรรมระหว่างดำเนินงาน

\subsection{ฝั่ง Front-End}

\begin{itemize}
  \item \textbf{โครงสร้างพื้นฐานพร้อมใช้งาน} Vite + React 18 + TypeScript + React Router 6 พร้อมชุดคำสั่ง \texttt{dev}, \texttt{build}, \texttt{lint}, \texttt{typecheck} ที่ใช้งานได้
  \item \textbf{ระบบออกแบบและเปลือกระบบ} Tailwind CSS พร้อมโทเคนสี แถบนำทางด้านข้าง แถบด้านบน และเมนูที่เปลี่ยนตามบทบาทผู้ใช้ทั้งสี่บทบาท
  \item \textbf{รองรับสองภาษาและสองธีม} ไทย–อังกฤษผ่าน i18next พร้อมโหมดสว่างและมืด ซึ่งเป็นงานที่แผนเดิมไม่ได้กำหนดไว้ใน Sprint 0
  \item \textbf{หน้าเข้าสู่ระบบที่ใช้งานได้จริง} รองรับการเข้าสู่ระบบสองช่องทาง คือผ่านอีเมล \texttt{@ku.ac.th} สำหรับผู้ยืม และบัญชีภายในสำหรับบทบาทอื่น พร้อมการคุ้มกันเส้นทางตามบทบาท
  \item \textbf{โครงหน้าจอครบทุกบทบาท} รวม 25 หน้า ครอบคลุมผู้ยืม เจ้าหน้าที่ ผู้อนุมัติ และผู้ดูแลระบบ
\end{itemize}

\noindent\textbf{ข้อจำกัดที่ต้องระบุให้ชัด} จาก 25 หน้าดังกล่าว มี \textbf{24 หน้าที่ยังเป็นหน้าเปล่า} (Placeholder) เนื้อหาจริงมีเพียงหน้าเข้าสู่ระบบหน้าเดียว และตัวดึงข้อมูลทั้งสองตัวที่มีอยู่คืนค่าจากข้อมูลจำลองที่เขียนไว้ในโค้ด ไม่ได้เรียก Back-End จริง

\subsection{ฝั่ง Back-End}

\begin{itemize}
  \item \textbf{โครงสร้างพื้นฐานพร้อมใช้งาน} NestJS พร้อม Prisma, \texttt{nestjs-trpc}, ตัวเชื่อม PostgreSQL และ Zod ประกอบเข้าด้วยกันใน \texttt{app.module.ts} เรียบร้อย
  \item \textbf{โครงสร้างฐานข้อมูลครบถ้วนและใช้งานได้จริง} \texttt{schema.prisma} นิยามตารางครบทั้ง 33 ตารางตรงตามการออกแบบรอบล่าสุด และมี Migration ที่สร้างตารางทั้ง 33 ตารางบนฐานข้อมูลจริงแล้วเมื่อวันที่ 5 สิงหาคม 2569
\end{itemize}

\noindent\textbf{ข้อจำกัดที่ต้องระบุให้ชัด} แม้จะลงทะเบียนโมดูล tRPC ไว้แล้ว แต่ \textbf{ยังไม่มี Procedure ใดถูกเขียนขึ้นเลยแม้แต่รายการเดียว} ฝั่ง Back-End จึงยังไม่มีจุดเชื่อมต่อใดที่ Front-End เรียกใช้ได้ นอกจากนี้ยังไม่มีสคริปต์สร้างข้อมูลตั้งต้น ทำให้ฐานข้อมูลที่สร้างขึ้นยังว่างเปล่า

\subsection{ฝั่งเอกสารและการออกแบบ}

\begin{itemize}
  \item \textbf{ข้อเสนอโครงการฉบับสมบูรณ์ 42 หน้า} ครอบคลุมบทบาทผู้ใช้ แผนภาพยูสเคส กระบวนการทำงาน การแบ่งระดับอุปกรณ์ ระเบียบการยืมและการจอง ระบบเครดิต ข้อมูลส่วนบุคคล การดำเนินงานของเจ้าหน้าที่ สถาปัตยกรรม และเทคโนโลยีที่ใช้
  \item \textbf{การออกแบบฐานข้อมูลรอบที่ 4} พร้อมพจนานุกรมข้อมูลครบทั้ง 33 ตาราง และแผนภาพความสัมพันธ์ของข้อมูล ซึ่งถูกนำไปเขียนเป็น Prisma schema จริงแล้ว
  \item \textbf{แผนภาพกระบวนการยืม–คืน} และ \textbf{แผนภาพยูสเคส} ที่วาดขึ้นใหม่ในเอกสาร
\end{itemize}

% =============================================================================
\section{การเปลี่ยนแปลงขอบเขตโครงการ}

\begin{tabularx}{\dimexpr\textwidth-\leftskip\relax}{@{} L{4.6cm} Y L{2.4cm} @{}}
\toprule
\hrow \thead{เรื่อง} & \thead{การเปลี่ยนแปลง} & \thead{ผลกระทบ} \\
\midrule
โครงสร้างฐานข้อมูล & ออกแบบใหม่ทั้งหมดเป็นรอบที่ 4 จาก 16 ตารางเป็น 33 ตาราง โดยใช้ \texttt{ResourceInfo} เป็นแกนกลางรวมอุปกรณ์กับห้อง & เป็นผลดี ทำให้การยืมอุปกรณ์และการจองสถานที่ใช้กระบวนการเดียวกัน \\
ตารางผูกชิ้นอุปกรณ์ & แผนเดิมกำหนดให้มีตาราง \texttt{ALLOCATION} สำหรับผูกการจองเข้ากับอุปกรณ์ชิ้นจริงตอนใกล้รับของ แต่การออกแบบรอบที่ 4 ไม่มีตารางนี้ & \textbf{ต้องแก้ไข} มิฉะนั้นงานจองล่วงหน้าใน Sprint 3 จะไม่มีที่เก็บข้อมูล \\
ตะกร้าข้ามหน่วยงาน & การออกแบบรอบที่ 4 ยังไม่มีสถานะ ``อนุมัติบางส่วน'' ซึ่งจำเป็นต่อการยืมของหลายภาคในคำขอเดียว & \textbf{ต้องแก้ไข} ก่อนถึง Sprint 3 \\
ฐานข้อมูล MongoDB & การออกแบบรอบที่ 4 ไม่มีข้อมูลชนิดรายการเหลืออยู่เลย เหตุผลที่เคยใช้อธิบายความจำเป็นของ MongoDB จึงใช้ไม่ได้อีก & ต้องตัดสินใจว่าจะคงไว้ในสถาปัตยกรรมหรือไม่ \\
\bottomrule
\end{tabularx}

% =============================================================================
\clearpage
\section{ปัญหาอุปสรรคและแนวทางแก้ไข}

\begin{tabularx}{\dimexpr\textwidth-\leftskip\relax}{@{} L{3.6cm} Y L{2.6cm} @{}}
\toprule
\hrow \thead{ปัญหา} & \thead{ผลกระทบและแนวทางแก้ไข} & \thead{ผู้รับผิดชอบ / กำหนด} \\
\midrule
สัญญา tRPC ยังไม่ถูกจัดทำ & \textbf{ผลกระทบ} เป็นต้นเหตุของงานที่ตกแผนแทบทั้งหมด เพราะทั้งสองฝั่งขาดจุดนัดหมายร่วมกัน จึงต่างคนต่างเขียนข้อมูลจำลองของตัวเอง และจะเจอปัญหาไม่ตรงกันเมื่อถึงเวลาเชื่อมต่อ \newline \textbf{แนวทาง} ให้กำหนดสัญญาเฉพาะวงจรหลักก่อน ประมาณ 6--8 Procedure ที่ครอบคลุมเข้าสู่ระบบ ดูรายการอุปกรณ์ สร้างคำขอ อนุมัติ ส่งมอบ และรับคืน โดยยังไม่ต้องครบทุกฟังก์ชัน & API + SA \newline ภายใน 9 ส.ค. \\
ฝั่ง Back-End ยังไม่มีฟังก์ชันงานจริง & \textbf{ผลกระทบ} เป้าหมายของ Sprint 1 ไม่มีทางสำเร็จตามกำหนดเดิม \newline \textbf{แนวทาง} ลดเป้าหมายเหลือเส้นทางเดียวที่ทะลุทุกชั้น คือ เข้าสู่ระบบ $\rightarrow$ ดูรายการอุปกรณ์ $\rightarrow$ สร้างคำขอ แล้วค่อยขยายใน Sprint 2 & BE-lead + BE-dev \newline ภายใน 13 ส.ค. \\
ไม่มีข้อมูลตั้งต้นในฐานข้อมูล & \textbf{ผลกระทบ} ต่อให้เขียน Procedure เสร็จก็ทดสอบไม่ได้ เพราะไม่มีข้อมูลให้เรียก \newline \textbf{แนวทาง} เขียนสคริปต์สร้างข้อมูลตั้งต้นสำหรับ 2 หน่วยงาน พร้อมผู้ใช้ครบทุกบทบาทและอุปกรณ์ตัวอย่าง & BE-dev \newline ภายใน 10 ส.ค. \\
งานทดสอบยังไม่เริ่มเลย & \textbf{ผลกระทบ} ยิ่งเริ่มช้า ยิ่งพบข้อบกพร่องช้า และค่าใช้จ่ายในการแก้ไขยิ่งสูง \newline \textbf{แนวทาง} เริ่มจากเกณฑ์การยอมรับของวงจรหลักก่อน ซึ่งเขียนได้ทันทีโดยไม่ต้องรอโค้ด & QA-lead \newline ภายใน 13 ส.ค. \\
ยังไม่มีระบบตรวจสอบอัตโนมัติ (CI) & \textbf{ผลกระทบ} โค้ดเสียถูกรวมเข้าสายหลักได้โดยไม่มีใครรู้ \newline \textbf{แนวทาง} ตั้ง GitHub Actions ให้รัน \texttt{lint} และ \texttt{typecheck} ทุกครั้งที่เปิดคำขอรวมโค้ด ซึ่งใช้เวลาตั้งค่าไม่เกินครึ่งวัน & FE-lead \newline ภายใน 10 ส.ค. \\
\bottomrule
\end{tabularx}

% =============================================================================
\section{การประเมินความเสี่ยง}

\begin{tabularx}{\dimexpr\textwidth-\leftskip\relax}{@{} L{4.2cm} L{1.5cm} L{1.7cm} Y @{}}
\toprule
\hrow \thead{ความเสี่ยง} & \thead{โอกาส} & \thead{ผลกระทบ} & \thead{แนวทางรับมือ} \\
\midrule
เชื่อมต่อ Front-End กับ Back-End ครั้งแรกแล้วพบว่าโครงสร้างข้อมูลไม่ตรงกัน & สูง & สูง & กำหนดสัญญา API ร่วมกันก่อนเขียนโค้ดเพิ่ม และให้ทั้งสองฝั่งสร้างชนิดข้อมูลจากสัญญาเดียวกัน \\
งานกระจุกตัวที่ Sprint 3 ซึ่งตรงกับช่วงสอบกลางภาค & สูง & สูง & ย้ายงานที่ไม่ขึ้นกับผู้อื่นขึ้นมาทำก่อนใน Sprint 2 และยืนยันว่าการจองสถานที่เป็นงานที่ตัดออกได้ \\
ฐานข้อมูลขาดตารางผูกชิ้นอุปกรณ์และสถานะอนุมัติบางส่วน & สูง & ปานกลาง & เพิ่มทั้งสองส่วนในการออกแบบรอบถัดไปก่อนเริ่ม Sprint 3 การแก้ตอนนี้ยังถูกกว่าแก้ตอนมีข้อมูลจริงแล้วมาก \\
กำลังคนฝั่ง Back-End ไม่พอเทียบสัดส่วนงาน & ปานกลาง & สูง & ย้ายสมาชิกที่ว่างจากฝั่ง Front-End มาช่วยเขียน Procedure ที่ไม่ซับซ้อน \\
งานที่ค้างสะสมถูกดันไปรวมกันที่ Sprint 2 & สูง & ปานกลาง & ยึดจุดตัดสินใจวันที่ 20 สิงหาคมอย่างเคร่งครัด หากตามไม่ทันให้ตัดงานที่ระบุว่าตัดออกได้ทันที ไม่ลากต่อไป Sprint ถัดไป \\
\bottomrule
\end{tabularx}

% =============================================================================
\clearpage
\section{แผนการดำเนินงานที่ปรับปรุงแล้ว}

\subsection{ช่วงที่เหลือของ Sprint 1 (8 – 13 สิงหาคม 2569)}

เนื่องจากเป้าหมายเดิมของ Sprint 1 ไม่มีทางสำเร็จภายในเวลาที่เหลือ คณะทำงานเสนอปรับเกณฑ์ ``เสร็จ'' ของ Sprint นี้ใหม่ดังนี้

\begin{tabularx}{\dimexpr\textwidth-\leftskip\relax}{@{} L{2.2cm} Y L{2.0cm} @{}}
\toprule
\hrow \thead{กำหนด} & \thead{งาน} & \thead{ผู้รับผิดชอบ} \\
\midrule
9 ส.ค. & กำหนดสัญญา tRPC ของวงจรหลัก 6--8 Procedure และตกลงร่วมกันทั้งสองฝั่ง & API + SA \\
10 ส.ค. & สคริปต์ข้อมูลตั้งต้น 2 หน่วยงาน พร้อมผู้ใช้ครบทุกบทบาท & BE-dev \\
10 ส.ค. & ตั้ง CI ให้รัน \texttt{lint} และ \texttt{typecheck} อัตโนมัติ & FE-lead \\
12 ส.ค. & Procedure สำหรับเข้าสู่ระบบและดูรายการอุปกรณ์ ที่เรียกใช้ได้จริง & BE-dev \\
13 ส.ค. & เปลี่ยนหน้ารายการอุปกรณ์จากข้อมูลจำลองเป็นข้อมูลจริงจาก Back-End & FE-dev \\
13 ส.ค. & เกณฑ์การยอมรับของวงจรหลัก & QA-lead \\
\bottomrule
\end{tabularx}

\vspace{0.8em}
\noindent\textbf{เกณฑ์ ``เสร็จ'' ที่ปรับใหม่} หน้ารายการอุปกรณ์แสดงข้อมูลที่ดึงมาจากฐานข้อมูลจริงผ่าน tRPC ได้สำเร็จ ถือเป็นการพิสูจน์ว่าทั้งสามชั้นของระบบต่อกันติด ส่วนงานสร้างคำขอ อนุมัติ ส่งมอบ และรับคืน ให้เลื่อนไป Sprint 2

\subsection{Sprint 2 (14 – 27 สิงหาคม 2569)}

\begin{itemize}
  \item งานที่เลื่อนมาจาก Sprint 1 ได้แก่ ยืนยันตัวตนจริง ควบคุมสิทธิ์ตามสมาชิกตามหน่วยงาน สร้างคำขอ อนุมัติ ส่งมอบ และรับคืน
  \item งานเดิมของ Sprint 2 ได้แก่ การอนุมัติตามระดับอุปกรณ์ การจัดเส้นทางอนุมัติตามหน่วยงานเจ้าของ การประเมินความเสียหาย และระบบเครดิต
  \item เพิ่มตารางผูกชิ้นอุปกรณ์และสถานะอนุมัติบางส่วนเข้าในโครงสร้างฐานข้อมูล
\end{itemize}

\noindent\textbf{ข้อควรระวัง} Sprint 2 จะแบกงานของสอง Sprint รวมกัน หากภายในวันที่ 20 สิงหาคมยังตามไม่ทัน คณะทำงานเสนอให้ตัดการจองสถานที่ (T3) และการส่งออกรายงานออกจากขอบเขตทันที ตามที่ระบุไว้แล้วว่าเป็นงานที่ตัดออกได้

\subsection{การประเมินกำหนดเสร็จของโครงการ}

จากการวิเคราะห์ข้างต้น คณะทำงานประเมินว่า \textbf{โครงการยังสามารถเสร็จได้ภายในวันที่ 1 ตุลาคม 2569 ตามกำหนดเดิม} โดยมีเงื่อนไขสองข้อ คือ ต้องจัดทำสัญญา API ให้เสร็จภายในสัปดาห์นี้ และต้องยอมตัดงานที่ระบุว่าตัดได้ออกจริงเมื่อถึงจุดตัดสินใจ หากเงื่อนไขข้อใดข้อหนึ่งไม่เป็นไปตามนี้ คณะทำงานจะรายงานขอเลื่อนกำหนดในรายงานความก้าวหน้าฉบับถัดไป

% =============================================================================
\section{ทรัพยากรและการสนับสนุนที่ต้องการเพิ่มเติม}

\begin{enumerate}
  \item \textbf{ฐานข้อมูล PostgreSQL สำหรับใช้ร่วมกันในทีม} ปัจจุบันแต่ละคนใช้ฐานข้อมูลบนเครื่องตัวเอง ทำให้ทดสอบร่วมกันไม่ได้ ค่าใช้จ่ายอยู่ในวงเงิน 1{,}500 บาทที่ประเมินไว้แล้ว
  \item \textbf{คำแนะนำจากอาจารย์ที่ปรึกษาเรื่องการคงไว้ซึ่ง MongoDB} เนื่องจากการออกแบบฐานข้อมูลรอบล่าสุดไม่มีข้อมูลที่จำเป็นต้องใช้ฐานข้อมูลแบบเอกสารเหลืออยู่ แต่การตัดออกจะทำให้ขอบเขตต่างจากที่เคยนำเสนอไว้
  \item \textbf{การยืนยันว่าการจองสถานที่ (T3) ตัดออกได้จริง} เพื่อให้คณะทำงานวางแผน Sprint 3 และ Sprint 4 ได้โดยไม่ต้องเผื่องานส่วนนี้
\end{enumerate}

% =============================================================================
\section{ความสำเร็จที่ควรบันทึกไว้}

\begin{itemize}
  \item \textbf{การออกแบบฐานข้อมูลเดินไปถึงขั้นใช้งานได้จริง} จากเอกสารออกแบบสู่ Prisma schema และ Migration ที่สร้างตารางครบ 33 ตารางบนฐานข้อมูลจริง เป็นงานที่เร็วกว่าแผน
  \item \textbf{เอกสารข้อเสนอโครงการมีความละเอียดสูง} 42 หน้าครอบคลุมกฎธุรกิจทั้งหมด ทำให้ทีมพัฒนามีแหล่งอ้างอิงเดียวที่ตรงกัน ลดการตีความผิดในภายหลัง
  \item \textbf{ฝั่ง Front-End ส่งมอบเกินแผนในบางส่วน} การรองรับสองภาษาและสองธีมไม่ได้อยู่ในแผน Sprint 0 แต่ทำเสร็จแล้ว ซึ่งเป็นงานที่หากทำทีหลังจะต้องแก้ทุกหน้าจอ
\end{itemize}

% =============================================================================
\clearpage
\section*{ภาคผนวก ก — ที่มาของตัวเลขความก้าวหน้า}
\addcontentsline{toc}{section}{ภาคผนวก ก — ที่มาของตัวเลขความก้าวหน้า}

ตัวเลขในหัวข้อ \ref{sec:dimension} ประเมินจากหลักฐานที่ตรวจสอบได้ใน Repository ณ วันที่ \reportdate\ ตามรายละเอียดต่อไปนี้

\begin{tabularx}{\dimexpr\textwidth-\leftskip\relax}{@{} L{3.4cm} L{1.2cm} Y @{}}
\toprule
\hrow \thead{มิติงาน} & \thead{ผลจริง} & \thead{หลักฐานที่ใช้ประเมิน} \\
\midrule
Back-End Design & 45\% & แบบจำลองข้อมูลประเมินที่ 85\% จาก \texttt{schema.prisma} ที่มีครบ 33 ตาราง และ Migration ที่สร้างตารางได้จริง 33 ตาราง หักส่วนที่ยังขาดคือตารางผูกชิ้นอุปกรณ์และสถานะอนุมัติบางส่วน ส่วนสัญญา API ประเมินที่ 5\% เพราะมีเพียงการลงทะเบียนโมดูล แต่ยังไม่มี Procedure ใดถูกเขียน \\
Front-End Design & 60\% & มีระบบออกแบบ เปลือกระบบ เมนูตามบทบาท และหน้าเข้าสู่ระบบครบ แต่ยังไม่มีแบบร่างหน้าจออีก 24 หน้า ส่วนการออกแบบกระบวนการทางธุรกิจถือว่าเสร็จแล้วจากเอกสารข้อเสนอโครงการ \\
Back-End Dev & 12\% & นับเฉพาะโครงสร้างพื้นฐานที่ประกอบเสร็จและ Migration ที่รันได้ ไม่มีไฟล์ใดใน \texttt{src/} ที่เป็นตรรกะทางธุรกิจ ไม่มี Router ไม่มีการยืนยันตัวตน และไม่มีสคริปต์ข้อมูลตั้งต้น \\
Front-End Dev & 20\% & จาก 25 หน้า มีเนื้อหาจริง 1 หน้า คิดเป็น 4\% บวกน้ำหนักของเปลือกระบบ การนำทาง การยืนยันตัวตนจำลอง และตัวดึงข้อมูล 2 ตัวที่โครงสร้างพร้อมสลับไปเรียก Back-End ได้ทันที \\
Test & 0\% & ไม่พบไฟล์ทดสอบที่เขียนขึ้นเอง มีเพียงไฟล์ตัวอย่างที่มากับโครงสร้างเริ่มต้นของ NestJS และไม่พบเอกสารแผนการทดสอบ \\
Documentation & 45\% & ข้อเสนอโครงการ 42 หน้า ครอบคลุมเนื้อหาของ Project Charter เกือบทั้งหมดและ SRS เป็นส่วนใหญ่ ส่วน SDS มีเฉพาะบทการออกแบบฐานข้อมูล ยังไม่มีคู่มือระบบและคู่มือผู้ใช้ \\
\bottomrule
\end{tabularx}

\vspace{1em}
\noindent\textbf{ข้อจำกัดของการประเมิน} ตัวเลขข้างต้นเป็นการประเมินจากปริมาณผลงานที่ปรากฏใน Repository ซึ่งเป็นหลักฐานที่ตรวจสอบย้อนกลับได้ทุกบรรทัด แต่การแปลงปริมาณผลงานเป็นร้อยละย่อมมีดุลพินิจของผู้ประเมินอยู่ด้วย โดยเฉพาะการกำหนดค่าน้ำหนักของแต่ละมิติงาน คณะทำงานจึงเลือกรายงานทั้งค่าน้ำหนักและที่มาของทุกตัวเลขไว้อย่างเปิดเผยในภาคผนวกนี้ เพื่อให้ผู้อ่านตรวจสอบสมมติฐานและปรับค่าตามที่เห็นสมควรได้

\end{document}
